\chapter{Experimental Methodology and Setup}
\label{ch:experiments}

\section{Benchmark Problem Instances}
To rigorously evaluate the scalability, convergence, and fairness of the optimization portfolio, three benchmark cohorts were generated reflecting the operational tiers of VJTI:
\begin{enumerate}
    \item \textbf{Small Instance (\texttt{small\_50x30})}: $N = 50$ students, $M = 30$ alumni mentors. Represents an intensive department-level pilot program (e.g., Computer Engineering departmental honors). This instance is computationally tractable for exact MILP branch-and-cut solvers, providing mathematical ground-truth bounds.
    \item \textbf{Medium Instance (\texttt{medium\_300x150})}: $N = 300$ students, $M = 150$ alumni mentors. Corresponds to a multi-departmental faculty cohort (e.g., all computing and electrical disciplines).
    \item \textbf{Large Instance (\texttt{large\_1000x500})}: $N = 1,000$ students, $M = 500$ alumni mentors. Represents the entire graduating undergraduate class of VJTI matched against the institutional alumni registry.
\end{enumerate}

\begin{table}[htbp]
    \centering
    \caption{Structural characteristics of synthetic benchmark instances.}
    \label{tab:instances}
    \begin{tabularx}{\textwidth}{lccccc}
        \toprule
        \textbf{Instance Identifier} & \textbf{Students ($N$)} & \textbf{Alumni ($M$)} & \textbf{Total Matches ($N \cdot K$)} & \textbf{Total Capacity ($\sum C_j$)} & \textbf{Degree ($K$)} \\
        \midrule
        \texttt{small\_50x30} & 50 & 30 & 100 & 120 & 2 \\
        \texttt{medium\_300x150} & 300 & 150 & 600 & 750 & 2 \\
        \texttt{large\_1000x500} & 1,000 & 500 & 2,000 & 2,500 & 2 \\
        \bottomrule
    \end{tabularx}
\end{table}

\section{Threshold Calibration Methodology ($\theta$ and $\delta$)}
A common flaw in academic assignment benchmarks is the arbitrary assignment of constraint thresholds, which frequently renders instances unintentionally infeasible. We developed an empirical, data-driven calibration procedure (\texttt{calibrate\_thresholds}):

\subsection{Rawlsian Quality Floor Calibration ($\theta$)}
The quality floor $\theta$ ensures that every student achieves viable mentorship. Setting $\theta$ too high causes mathematical infeasibility; setting it too low renders the constraint inactive. We compute $\theta$ as a calibrated fraction of the unconstrained Greedy baseline's worst-student quality $Z_2^{\text{greedy}}$:
\begin{equation}
    \theta = \alpha_{\theta} \cdot Z_2^{\text{greedy}}, \quad \alpha_{\theta} = 0.85
\end{equation}
Because $\alpha_{\theta} < 1.0$, the greedy baseline itself is guaranteed to be feasible under the calibrated $\theta$, ensuring a non-empty feasible set.

\subsection{Departmental Parity Cap Calibration ($\delta$)}
The branch parity parameter $\delta$ bounds the allowable difference between the most-advantaged and least-advantaged departmental averages. We derive $\delta$ from the empirical branch spread of the unconstrained Greedy assignment:
\begin{equation}
    \Delta_{\text{greedy}} = \max_{b \in \mathcal{B}} \overline{Q}_b^{\text{greedy}} - \min_{b \in \mathcal{B}} \overline{Q}_b^{\text{greedy}}
\end{equation}
\begin{equation}
    \delta = \alpha_{\delta} \cdot \Delta_{\text{greedy}}, \quad \alpha_{\delta} = 1.20
\end{equation}
Providing a 20\% slack margin above the unconstrained spread ensures feasibility while strictly penalizing extreme departmental divergence.

\begin{table}[htbp]
    \centering
    \caption{Calibrated institutional thresholds across benchmark cohorts.}
    \label{tab:calibration}
    \begin{tabularx}{\textwidth}{lccccc}
        \toprule
        \textbf{Cohort Instance} & \textbf{Greedy $Z_2$} & \textbf{Calibrated $\theta$ ($\alpha=0.85$)} & \textbf{Greedy Spread $\Delta$} & \textbf{Calibrated $\delta$ ($\alpha=1.20$)} & \textbf{Feasible?} \\
        \midrule
        \texttt{small\_50x30} & 1.0230 & \textbf{0.8696} & 0.3258 & \textbf{0.3910} & \textbf{Yes} \\
        \texttt{medium\_300x150} & 1.0434 & \textbf{0.8869} & 0.3295 & \textbf{0.3953} & \textbf{Yes} \\
        \texttt{large\_1000x500} & 1.3211 & \textbf{1.1229} & 0.2386 & \textbf{0.2863} & \textbf{Yes} \\
        \bottomrule
    \end{tabularx}
\end{table}

\section{Evaluation Budget Equalization}
To guarantee rigorous and unbiased scientific comparisons between meta-heuristics, all evolutionary algorithms (NSGA-II, SPEA2, MOEA/D, and Hybrid Memetic) were configured under an strictly equal evaluation budget:
\begin{equation}
    \text{Budget} \approx N_{\text{pop}} \times N_{\text{gen}} \approx 1,600 \text{ function evaluations (on \texttt{small\_50x30})}
\end{equation}
\begin{itemize}
    \item \textbf{NSGA-II}: Population size $N_{\text{pop}} = 40$, Generations $N_{\text{gen}} = 40$ (1,600 evals).
    \item \textbf{SPEA2}: Population size $N_{\text{pop}} = 40$, Generations $N_{\text{gen}} = 40$ (1,600 evals).
    \item \textbf{MOEA/D}: 6 partitions, 57 generations ($\approx 1,600$ evals).
    \item \textbf{Hybrid Memetic}: Population size $N_{\text{pop}} = 40$, Generations $N_{\text{gen}} = 35$, plus local search evaluation accounting ($\approx 1,600$ total function evaluations).
\end{itemize}

\section{Unified Shared Hypervolume Reference Setup}
A critical methodological vulnerability identified in prior literature is the use of per-method fallback reference points (e.g., multiplying each algorithm's observed maximum by 1.1). When algorithms achieve disparate maximums, comparing their raw hypervolume values is mathematically invalid.

To eliminate this flaw, we implemented a \textbf{Unified Shared Hypervolume Reference System}:
\begin{enumerate}
    \item \textbf{Front Pooling}: For each instance, we gather the non-dominated Pareto objective vectors from all eight methods across all 30 random seeds, combined with the exact MILP Pareto front on the small instance:
    \begin{equation}
        \mathcal{F}_{\text{all}} = \bigcup_{m \in \text{Methods}} \bigcup_{s=1}^{30} \mathcal{F}_{m,s} \cup \mathcal{F}_{\text{MILP}}
    \end{equation}
    \item \textbf{Global Ideal and Nadir Computation}: We compute the global ideal point $\mathbf{z}^*$ and nadir point $\mathbf{z}^{\text{nad}}$ across all objectives in minimization space ($[-Z_1, -Z_2, Z_3]$):
    \begin{equation}
        z_k^* = \min_{\mathbf{f} \in \mathcal{F}_{\text{all}}} f_k, \quad z_k^{\text{nad}} = \max_{\mathbf{f} \in \mathcal{F}_{\text{all}}} f_k, \quad \forall k \in \{1, 2, 3\}
    \end{equation}
    \item \textbf{Normalized Objective Space}: Every candidate objective vector $\mathbf{f}$ is normalized to the unit hypercube:
    \begin{equation}
        \widetilde{f}_k = \frac{f_k - z_k^*}{z_k^{\text{nad}} - z_k^* + \epsilon}, \quad \widetilde{\mathbf{f}} \in [0, 1]^3
    \end{equation}
    \item \textbf{Fixed Invariant Reference Point}: Hypervolume is computed against a strictly fixed reference point:
    \begin{equation}
        \mathbf{r} = [1.1, 1.1, 1.1]
    \end{equation}
\end{enumerate}
Under this formulation, two identical Pareto fronts yield mathematically identical HV scores, and a strictly dominated front is guaranteed to receive a lower HV.

\section{Multi-Seed Protocol and Statistical Significance Testing}
All stochastic algorithms were executed across $S = 30$ independent random seeds ($s \in \{42, \dots, 71\}$). For deterministic baselines (Greedy, MCF), evaluations run once and are replicated across seed rows for paired statistical integrity.

To evaluate whether the Hybrid Memetic algorithm provides statistically significant improvements over competing baselines, we conduct paired two-sided non-parametric \textbf{Wilcoxon signed-rank tests}. We report the test statistic $T$, the asymptotic $p$-value, and the rank-biserial effect size $r$:
\begin{equation}
    r = \frac{T^+ - T^-}{T^+ + T^-} = 1 - \frac{2T}{\frac{n(n+1)}{2}}
\end{equation}
where an effect size $r \ge 0.5$ denotes a large operational advantage.
